%
% reversi_270_iperparametri.tex
%
% This file is part of the reversi program
% http://github.com/rcrr/reversi-v2
%
% Author Roberto Corradini mailto:rob_corradini@yahoo.it
% Copyright 2026 Roberto Corradini. All rights reserved.
%
% License
%
% This program is free software; you can redistribute it and/or modify it
% under the terms of the GNU General Public License as published by the
% Free Software Foundation; either version 3, or (at your option) any
% later version.
%
% This program is distributed in the hope that it will be useful,
% but WITHOUT ANY WARRANTY; without even the implied warranty of
% MERCHANTABILITY or FITNESS FOR A PARTICULAR PURPOSE. See the
% GNU General Public License for more details.
%
% You should have received a copy of the GNU General Public License
% along with this program; if not, write to the Free Software
% Foundation, Inc., 59 Temple Place - Suite 330, Boston, MA  02111-1307, USA
% or visit the site <http://www.gnu.org/licenses/>.
%

%-----------------------------------------------------------------------------
% Capitolo: Gli hyper-parameter.
%
% Tredicesima fase: lo sweep della regolarizzazione ridge sul campione
% (0104, identico a 0103 salvo il parametro), con la campagna a passo
% unitario da 0 a 10: il minimo e' la regione piatta [0, 2], la curva
% sale in modo monotono da 3 a 10 (7.10); la normalizzazione
% alpha = lambda / M spinge il ginocchio della curva verso l'alto a
% questa scala del dato; il valore in uso (0.01) e' al centro del minimo.

\chapter{Gli hyper-parameter}

\section{La regolarizzazione ridge: la curva del minimo}

Eseguiamo una campagna di messa a punto degli hyper-parameter del
modello logistico, usiamo il miglior modello costruito nei capitoli
precedenti, con i dieci pattern più OptMob, e il batch di training da
14 milioni di posizioni. Incominciamo dalla regolarizzazione di
Ridge, sempre fissa a 0.01 in tutti i modelli sviluppati fino ad ora,
la domanda è se esista un valore che abbassi la RMSE di validazione
del campione.

La scansione gira su un unico modello nuovo (0104), identico al
campione 0103 salvo il parametro: prima la metà (0.005), poi il doppio
(0.02), poi piccoli passi nella direzione del miglioramento (0.04,
0.03, 0.015 a sondare il centro, 0.001, 0.025, 0.05, 0.1 a
esplorare l'intervallo largo, 0.0 e 0.5 a chiudere i due capi, e
infine la campagna completa a passo unitario, da 0 a 10).

La tabella~\ref{tab:ridge-curva} riporta i punti della campagna a
passo unitario.

\begin{table}[ht]
\centering
\begin{tabular}{lrr}
\toprule
ridge & RMSE validazione & iterazioni \\
\midrule
0  & 7.05 & 301 \\
1  & 7.05 & 261 \\
2  & 7.05 & 248 \\
3  & 7.06 & 237 \\
4  & 7.06 & 219 \\
5  & 7.07 & 204 \\
6  & 7.08 & 193 \\
7  & 7.08 & 198 \\
8  & 7.09 & 197 \\
9  & 7.10 & 194 \\
10 & 7.10 & 191 \\
\bottomrule
\end{tabular}
\caption{La RMSE di validazione al variare della regolarizzazione
ridge, a passo unitario da 0 a 10 (il modello 0104, identico a 0103
salvo il parametro): la RMSE di validazione e le iterazioni di
ottimizzazione.}
\label{tab:ridge-curva}
\end{table}

La campagna a passo unitario chiude il quadro. Il minimo è la regione
piatta da 0 a 2 — 7.05 ai punti interi, 7.04 al miglior punto
osservato, 0.01 — e da 3 in poi la curva sale in modo monotono, di un
centesimo alla volta, fino a 7.10 a 10: non c'è alcun minimo nascosto
nel tratto, la funzione è piatta sul primo intervallo e regolarmente
crescente sul secondo. Il valore in uso — 0.01 — sta al centro del
minimo largo, e il bordo destro, netto e regolare, è la firma di una
penalità che comincia a contare.

Perché il pianoro è così largo? La risposta sta nella
normalizzazione: la pipeline scala la penalità per il numero di
posizioni, $\alpha = \lambda / M$, in modo che i criteri di
convergenza non dipendano dalla dimensione del corpus. Su
quattordici milioni di righe, un $\lambda$ nominale fisso pesa dunque
molto meno per campione che su un corpus piccolo, e il ginocchio della
curva viene spinto di conseguenza verso l'alto. Il ritocco della
regolarizzazione resta privo di beneficio: il valore in uso è al
centro del minimo, e la salita a 10 conferma che oltre quel piano il
parametro comincia a costare.
